\SixthArticle{41}

Setzen wir nun $P=x^\mu P'$, sodass sich ergibt
\[
\begin{aligned}
\frac{\mathrm{d}P}{\mathrm{d}x}&=\mu x^{\mu-1}P'+x^\mu\frac{\mathrm{d}P'}{\mathrm{d}x}\\
\frac{\mathrm{dd}P}{\mathrm{d}x^2}&=(\mu\mu-\mu)x^{\mu-2}P'+2\mu x^{\mu-1}\frac{\mathrm{d}P'}{\mathrm{d}x}+x^\mu\frac{\mathrm{dd}P'}{\mathrm{d}x^2}
\end{aligned}
\]
\SourceLeaf{219}
Setzt man diese Werte in 80 ein und dividiert durch $x^{\mu-1}$, so erhält man
\[
\begin{aligned}
0={}&P'\{\alpha\beta x-(\gamma-(\alpha+\beta+1)x)\mu-(1-x)(\mu\mu-\mu)\}\\
&-\frac{\mathrm{d}P'}{\mathrm{d}x}\{\gamma-(\alpha+\beta+1)x+2\mu(1-x)\}x\\
&-\frac{\mathrm{dd}P'}{\mathrm{d}x^2}(xx-x^3)
\end{aligned}
\]
Der Faktor von $P'$ in dieser Formel wird durch $x$ teilbar, wenn man $\mu=0$ oder $\mu=1-\gamma$ setzt. Der letztere Wert liefert
\[
\begin{aligned}
0={}&P'(\alpha\beta+\alpha+\beta+1-2\gamma-\alpha\gamma-\beta\gamma+\gamma\gamma)\\
&-\frac{\mathrm{d}P'}{\mathrm{d}x}(2-\gamma-(\alpha+\beta+3-2\gamma)x)\\
&-\frac{\mathrm{dd}P'}{\mathrm{d}x^2}(x-xx)
\end{aligned}
\]
Der Vergleich dieser Gleichung mit 80, deren Form genau ähnlich ist, zeigt, dass das, was dort $P$, $\alpha$, $\beta$, $\gamma$ war, hier $P'$, $\alpha+1-\gamma$, $\beta+1-\gamma$, $2-\gamma$ ist. Da wir die allgemeine Lösung jener Gleichung angegeben haben, ist ersichtlich, dass für $P'$ die folgende Formel gelten wird:
\[
\begin{aligned}
P'={}&MF(\alpha+1-\gamma,\beta+1-\gamma,2-\gamma,x)\\
&+NF(\alpha+1-\gamma,\beta+1-\gamma,\alpha+\beta+1-\gamma,1-x)
\end{aligned}
\]
wobei $M$, $N$ konstante Größen bezeichnen, oder
\[
\begin{aligned}
&\text{[83]}\qquad F(\alpha,\beta,\gamma,x)\\&=Mx^{1-\gamma}F(\alpha+1-\gamma,\beta+1-\gamma,2-\gamma,x)\\
&+Nx^{1-\gamma}F(\alpha+1-\gamma,\beta+1-\gamma,\alpha+\beta+1-\gamma,1-x)
\end{aligned}
\]
wobei die Konstanten $M$, $N$ von den Elementen $\alpha$, $\beta$, $\gamma$ abhängen werden.

\SixthArticle{42}

Aus Gleichung 82 folgt
\[
\begin{aligned}
&F(\alpha+1-\gamma,\beta+1-\gamma,2-\gamma,x)\\&=(1-x)^{\gamma-\alpha-\beta}F(1-\alpha,1-\beta,2-\gamma,x)\\[5pt]
&F(\alpha+1-\gamma,\beta+1-\gamma,\alpha+\beta+1-\gamma,1-x)\\&=x^{\gamma-1}F(\alpha,\beta,\alpha+\beta+1-\gamma,1-x)
\end{aligned}
\]
Setzt man daher
\[
\frac{1}{N}=f(\alpha,\beta,\gamma),\qquad-\frac{M}{N}=g(\alpha,\beta,\gamma)
\]
\SourceLeaf{220}
so wird Gleichung 83 zu
\[
\begin{aligned}
&F(\alpha,\beta,\alpha+\beta+1-\gamma,1-x)\\&=f(\alpha,\beta,\gamma)F(\alpha,\beta,\gamma,x)\\
&\quad+g(\alpha,\beta,\gamma)(1-x)^{\gamma-\alpha-\beta}x^{1-\gamma}F(1-\alpha,1-\beta,2-\gamma,x)
\end{aligned}
\]
Mithilfe von Formel 82 lässt sich dieselbe Gleichung auch folgendermaßen schreiben:
\[
\begin{aligned}
&x^{1-\gamma}F(\alpha+1-\gamma,\beta+1-\gamma,\alpha+\beta+1-\gamma,1-x)=\\
&f(\alpha,\beta.\gamma)(1-x)^{\gamma-\alpha-\beta}F(\gamma-\alpha,\gamma-\beta,\gamma,x)
\\&+g(\alpha,\beta,\gamma)x^{1-\gamma}F(\alpha+1-\gamma,\beta+1-\gamma,2-\gamma,x)
\end{aligned}
\]
Oder man dividiert durch \(x^{1-\gamma}\) und ersetzt \(\alpha,\beta,\gamma\) jeweils durch \(\alpha+1-\gamma,\beta+1-\gamma,2-\gamma\); dann erhält man
\[
\begin{aligned}
&F(\alpha,\beta,\alpha+\beta+1-\gamma,1-x)\\
&\quad=g(\alpha+1-\gamma,\beta+1-\gamma,2-\gamma)F(\alpha,\beta,\gamma,x)\\
&\qquad+f(\alpha+1-\gamma,\beta+1-\gamma,2-\gamma)\\&\qquad(1-x)^{\gamma-\alpha-\beta}x^{1-\gamma}F(1-\alpha,1-\beta,2-\gamma,x)
\end{aligned}
\]
Da dies mit der vorhergehenden Formel identisch sein muss, haben wir
\[
g(\alpha,\beta,\gamma)=f(\alpha+1-\gamma,\beta+1-\gamma,2-\gamma)
\]
und somit
\[
\begin{aligned}
\text{[84]}\quad &F(\alpha,\beta,\alpha+\beta+1-\gamma,1-x)\\
&=f(\alpha,\beta,\gamma)F(\alpha,\beta,\gamma,x)\\
&\quad+f(\alpha+1-\gamma,\beta+1-\gamma,2-\gamma)\\&\qquad(1-x)^{\gamma-\alpha-\beta}x^{1-\gamma}F(1-\alpha,1-\beta,2-\gamma,x)
\end{aligned}
\]
\SixthArticle{43}
Um nun die Beschaffenheit der Funktion \(f(\alpha,\beta,\gamma)\) zu ermitteln, setzen wir \(x=0\). Dann ist offenbar \(F(\alpha,\beta,\gamma,x)=1.\ x^{1-\gamma}=0\), sofern \(1-\gamma\) eine positive Größe ist. Nach Gleichung 48 haben wir aber
\[
F(\alpha,\beta,\alpha+\beta+1-\gamma,1)=\frac{\Pi(\alpha+\beta-\gamma)\Pi(-\gamma)}{\Pi(\alpha-\gamma)\Pi(\beta-\gamma)}
\]
Unter derselben Einschränkung ist also bewiesen:
\[
\text{[85]}\qquad f(\alpha,\beta,\gamma)=\frac{\Pi(\alpha+\beta-\gamma)\Pi(-\gamma)}{\Pi(\alpha-\gamma)\Pi(\beta-\gamma)}
\]
Dass diese Formel allgemein gilt, beweisen wir jedoch folgendermaßen.
\par\SourceLeaf{221}
Durch Differenzieren der Gleichung 84 erhält man
\[
\begin{aligned}
&-\frac{\alpha\beta}{\alpha+\beta+1-\gamma}F(\alpha+1,\beta+1,\alpha+\beta+2-\gamma,1-x)\\
&=\frac{\alpha\beta}{\gamma}f(\alpha,\beta,\gamma)F(\alpha+1,\beta+1,\gamma+1.\ x)\\
&\quad+f(\alpha+1-\gamma,\beta+1-\gamma,2-\gamma)(1-x)^{\gamma-\alpha-\beta-1}x^{-\gamma}\\&\qquad\Bigl\{((1-\gamma)(1-x)-(\gamma-\alpha-\beta)x)F(1-\alpha,1-\beta,2-\gamma,x)\\
&\qquad+\frac{(1-\alpha)(1-\beta)}{2-\gamma}(x-xx)F(2-\alpha,2-\beta,3-\gamma,x)\Bigr\}
\end{aligned}
\]
Nach Formel IX aus Artikel 10 ergibt sich aber, wenn man \(\alpha,\beta,\gamma\) durch \(-\alpha,-\beta,1-\gamma\) ersetzt,
\[
\begin{aligned}
&(1-\gamma)(2-\gamma)F(-\alpha,-\beta,1-\gamma,x)\\&=(2-\gamma)(1-\gamma+(\alpha+\beta-1)x)F(1-\alpha,1-\beta,2-\gamma,x)\\
&\quad+(1-\alpha)(1-\beta)(x-xx)F(2-\alpha,2-\beta,3-\gamma,x)
\end{aligned}
\]
Dadurch geht die vorhergehende Gleichung über in
\[
\begin{aligned}
&F(\alpha+1,\beta+1,\alpha+\beta+2-\gamma,1-x)\\
&=-\frac{\alpha+\beta+1-\gamma}{\gamma}f(\alpha,\beta,\gamma)F(\alpha+1,\beta+1,\gamma+1.\ x)\\
&\quad-\frac{(\alpha+\beta+1-\gamma)(1-\gamma)}{\alpha\beta}f(\alpha+1-\gamma,\beta+1-\gamma,2-\gamma)\\&\qquad(1-x)^{\gamma-\alpha-\beta-1}x^{-\gamma}F(-\alpha,-\beta,1-\gamma,x)
\end{aligned}
\]
Ersetzt man dagegen in Gleichung 84 \(\alpha,\beta,\gamma\) durch \(\alpha+1,\beta+1,\gamma+1\), so folgt
\[
\begin{aligned}
&F(\alpha+1,\beta+1,\alpha+\beta+2-\gamma,1-x)\\
&\quad=f(\alpha+1,\beta+1,\gamma+1)F(\alpha+1,\beta+1,\gamma+1,x)\\
&\qquad+f(\alpha+1-\gamma,\beta+1-\gamma,1-\gamma)\\&\qquad(1-x)^{\gamma-\alpha-\beta-1}x^{-\gamma}F(-\alpha,-\beta,1-\gamma,x)
\end{aligned}
\]
Da leicht zu erkennen ist, dass diese beiden Gleichungen identisch sein müssen, gilt allgemein
\[
f(\alpha+1,\beta+1,\gamma+1)=\frac{\alpha+\beta+1-\gamma}{-\gamma}f(\alpha,\beta,\gamma)
\]
oder, wenn man \(\alpha,\beta,\gamma\) durch \(\alpha-1,\beta-1,\gamma-1\) ersetzt,
\[
\begin{aligned}
f(\alpha,\beta,\gamma)&=\frac{\alpha+\beta-\gamma}{1-\gamma}f(\alpha-1,\beta-1,\gamma-1)\\
&=\frac{\alpha+\beta-\gamma\,.\,\alpha+\beta-\gamma-1}{1-\gamma\,.\,2-\gamma}f(\alpha-2,\beta-2,\gamma-2)
\end{aligned}
\]
usw. Daraus lässt sich leicht schließen, dass allgemein für jeden ganzzahligen Wert von \(k\)
\[
f(\alpha,\beta,\gamma)=\frac{\Pi(\alpha+\beta-\gamma)\Pi(-\gamma)}{\Pi(\alpha+\beta-\gamma-k)\Pi(k-\gamma)}.f(\alpha-k,\beta-k,\gamma-k)
\]
gilt.
\par\SourceLeaf{222}
Sofern aber $1-(\gamma-k)$ oder $k+1-\gamma$ eine positive Größe ist, haben wir bewiesen (Formel 85):
\[
f(\alpha-k,\beta-k,\gamma-k)
=\frac{\Pi(\alpha+\beta-\gamma-k)\Pi(k-\gamma)}{\Pi(\alpha-\gamma)\Pi(\beta-\gamma)}
\]
Da man nun $k$, welchen Wert auch $\gamma$ haben mag, stets so groß wählen kann, dass $k+1-\gamma$ eine positive Größe wird, gilt \textit{allgemein}
\[
f(\alpha,\beta,\gamma)
=\frac{\Pi(\alpha+\beta-\gamma)\Pi(-\gamma)}{\Pi(\alpha-\gamma)\Pi(\beta-\gamma)}
\]
und daher
\[
f(\alpha+1-\gamma,\beta+1-\gamma,2-\gamma)
=\frac{\Pi(\alpha+\beta-\gamma)\Pi(\gamma-2)}{\Pi(\alpha-1)\Pi(\beta-1)}
\]
sodass unsere Formel lautet
\[
\begin{aligned}[c]
[86]\quad &F(\alpha,\beta,\alpha+\beta+1-\gamma,1-x)\\
&=\frac{\Pi(\alpha+\beta-\gamma)\Pi(-\gamma)}{\Pi(\alpha-\gamma)\Pi(\beta-\gamma)}F(\alpha,\beta,\gamma,x)\\
&\quad+\frac{\Pi(\alpha+\beta-\gamma)\Pi(\gamma-2)}{\Pi(\alpha-1)\Pi(\beta-1)}x^{1-\gamma}\\&\qquad(1-x)^{\gamma-\alpha-\beta}F(1-\alpha,1-\beta,2-\gamma,x)
\end{aligned}
\]
oder, wenn man $\gamma$ durch $\alpha+\beta+1-\gamma$ ersetzt,
\[
\begin{aligned}[c]
[87]\quad &F(\alpha,\beta,\gamma,1-x)\\
&=\frac{\Pi(\gamma-1)\Pi(\gamma-\alpha-\beta-1)}{\Pi(\gamma-\alpha-1)\Pi(\gamma-\beta-1)}F(\alpha,\beta,\alpha+\beta+1-\gamma,x)\\
&\quad+\frac{\Pi(\gamma-1)\Pi(\alpha+\beta-\gamma-1)}{\Pi(\alpha-1)\Pi(\beta-1)}x^{\gamma-\alpha-\beta}\\&\qquad(1-x)^{1-\gamma}F(1-\alpha,1-\beta,\gamma+1-\alpha-\beta,x)
\end{aligned}
\]

Wenn man es vorzieht, kann man schreiben

in Formel 86
\[
\begin{aligned}&\text{für }(1-x)^{\gamma-\alpha-\beta}F(1-\alpha,1-\beta,2-\gamma,x)\\&\ .\ .\ .\ .\ F(\alpha+1-\gamma,\beta+1-\gamma,2-\gamma,x)\end{aligned}
\]
in Formel 87
\[
\begin{aligned}&\text{für }(1-x)^{1-\gamma}F(1-\alpha,1-\beta,\gamma+1-\alpha-\beta,x)\\&\ .\ .\ .\ .\ F(\gamma-\alpha,\gamma-\beta,\gamma+1-\alpha-\beta,x)\end{aligned}
\]

\SixthArticle{44}

Wird also dem vierten Element einer Reihe unserer Form ein Wert zwischen $0{,}5$ und $1$ gegeben, so lässt sich ihrer langsameren Konvergenz durch die vorhergehenden Formeln abhelfen. Diese zerlegen sie nämlich in zwei andere ähnliche Reihen,
\SourceLeaf{223}
die umso schneller konvergieren, je langsamer die ursprüngliche Reihe konvergierte. Auszunehmen sind jedoch die besonderen Fälle, in denen diese Umformung nicht gelingt, nämlich immer dann, wenn in der umzuformenden Reihe die Differenz zwischen dem dritten Element und der Summe der beiden ersten Elemente eine ganze Zahl ist. Ist nämlich in Formel 86 $\gamma$ $=0$ oder gleich einer negativen ganzen Zahl, so wird $F(\alpha,\beta,\gamma,x)$ offenbar eine unbrauchbare Reihe (Artikel 2) und der Faktor $\Pi(\gamma-2)$ unendlich. Ist dagegen $\gamma$ eine positive ganze Zahl größer als eins, so werden $F(1-\alpha,1-\beta,2-\gamma,x)$ und $F(\alpha+1-\gamma,\beta+1-\gamma,2-\gamma,x)$ unbrauchbare Reihen und $\Pi(-\gamma)$ unendlich. Ist schließlich $\gamma=1$, so leiden die beiden umgeformten Reihen $F(\alpha,\beta,\gamma,x)$ und $F(1-\alpha,1-\beta,2-\gamma,x)$ oder $F(\alpha+1-\gamma,\beta+1-\gamma,2-\gamma,x)$, die folglich mit $F(\alpha,\beta,\gamma,x)$ identisch wird, zwar nicht unter diesem Nachteil. Dennoch ist die Umformung nutzlos, da jede der beiden umgeformten Reihen mit dem unendlichen Koeffizienten $\Pi(-1)$ multipliziert wird. Es wird sich daher lohnen zu zeigen, wie sich auch in diesen Fällen eine langsamere Konvergenz in eine schnellere überführen lässt.

